%%%%%%%%%%%%%%%%%%%%%%%%%%%%%%%%%%%%%%%%%%%%%%%%%%%%%%%%%%%%%%%%%%%%%%%%%%%%%%
%% Example Ph.D. thesis using tvthesis class
%%
%% The PhD thesis uses a custom book format (17x24cm).
%% To embed it in A4 paper with crop marks, use:
%%   \documentclass[phd,english,a4paper,showcrop]{tvthesis}
%%
%% Compile with: pdflatex example-phd.tex (twice)
%%%%%%%%%%%%%%%%%%%%%%%%%%%%%%%%%%%%%%%%%%%%%%%%%%%%%%%%%%%%%%%%%%%%%%%%%%%%%%
\documentclass[phd,english,lof,lot]{tvthesis}

\doctorate{Computer Science}
\cycle{XXXVIII}

\title{Techniques for Defensive Programming}
\subtitle{Design, Implementation, and Evaluation}
\author{John Doe}

\advisor{Prof.~Alice Smith}
\committeemember{Prof.~Bob Johnson}
\committeemember{Prof.~Carol Williams}

\reviewer{Prof.~David Brown}
\reviewer{Prof.~Emily Davis}

\AcademicYear{A.Y. 2025/2026}

\authoremail{john.doe@example.com}
\authoraddress{Department of Civil Engineering and Computer Engineering\\
  Universit\`a degli Studi di Roma Tor Vergata\\
  Via del Politecnico 1, 00133 Roma, Italy}
\website{https://www.example.com/jdoe}

\usepackage{lipsum}

\begin{document}

\maketitle

\frontmatter

\dedication{To my family,\\who supported me throughout this journey}

\tableofcontents


\begin{abstract}
\lipsum[1]
\end{abstract}


\begin{acknowledgments}
\lipsum[2]
\end{acknowledgments}



\mainmatter


\chapter*{Abstract}
\label{abstract}
\addcontentsline{toc}{chapter}{Abstract}
\markboth{\textsc{Abstract}}{\textsc{Abstract}}

\lipsum[3][3-4]

\section*{Chapter Organization}

In Chapter~\ref{ch:introduction}, we introduce the research problems.
Chapter~\ref{ch:context} frames the research context.
Chapter~\ref{ch:conclusions} draws conclusions.



\chapter{Introduction}
\label{ch:introduction}

Over the past decades, the increasing availability of multi-core processors
has fundamentally changed the landscape of high-performance
computing~\cite{sutter2005}. 

In this thesis, we explore the simulation research field to provide
methodological and technical solutions towards the generation of efficient
parallel simulation models from sequential ones.


\section{Research Context}

\lipsum[4]

\section{Contributions}

The main contributions of this thesis are:
\begin{enumerate}
    \item \lipsum[5][1]
    \item \lipsum[5][2]
    \item \lipsum[5][3]
\end{enumerate}

\lipsum[6]



\chapter{Research Context and Achieved Results}
\label{ch:context}

\lipsum[7]


\section{More Lorem Ipsum}

\lipsum[8-9]

\section{Consectetuer}

\lipsum[10]

Figure~\ref{fig:speedup} illustrates the theoretical speedup as a function
of the number of processing units for different parallelizable fractions.

\begin{figure}[t]
    \centering
    \includegraphics[width=0.75\textwidth]{example-image}
    \caption{Theoretical speedup according to Amdahl's Law for varying
    fractions $P$.}
    \label{fig:speedup}
\end{figure}

Table~\ref{tab:benchmarks} summarises the benchmark configurations used
in our experimental evaluation.

\begin{table}[t]
    \centering
    \caption{Benchmark configurations for the experimental evaluation.}
    \label{tab:benchmarks}
    \begin{tabular}{lrrr}
        \toprule
        \textbf{Benchmark} & \textbf{Cores} & \textbf{Workload} & \textbf{Think Time} \\
        \midrule
        A & 2   & 500  & 1.0 \\
        B & 4   & 200  & 0.5 \\
        C & 16 & 100 & 0.1 \\
        \bottomrule
    \end{tabular}
\end{table}

\lipsum[11-12]



\chapter{Conclusions}
\label{ch:conclusions}

\lipsum[13-14]



\begin{thebibliography}{9}

\bibitem{sutter2005}
H.~Sutter.
\newblock The free lunch is over: A fundamental turn toward concurrency
  in software.
\newblock {\em Dr.\ Dobb's Journal}, 30(3):202--210, 2005.

\end{thebibliography}

\end{document}
